\section{Domain Analysis}
\subsection{Application Context and Utilized Technologies}

The primary objective of this laboratory work is to implement a \textbf{binary signal
acquisition and conditioning system} for a digital temperature sensor using
\textbf{FreeRTOS} on the ESP32 microcontroller. The application demonstrates a
complete sensor pipeline: periodic raw data acquisition, hysteresis-based threshold
detection with debounce filtering, and structured reporting on both an LCD display and
the serial interface.

The system uses the \textbf{DS18B20} one-wire digital temperature sensor. The sensor
communicates over the \textbf{OneWire} protocol on a single GPIO pin (GPIO\,4), and is
read through the \textit{DallasTemperature} library in a non-blocking configuration
(9-bit resolution, \texttt{setWaitForConversion(false)}). This pattern issues the
temperature conversion request in one acquisition cycle and retrieves the result in the
next, avoiding blocking delays inside the FreeRTOS task.

The signal conditioning stage implements two classical digital signal processing
techniques:
\begin{itemize}
  \item \textbf{Hysteresis:} The alert turns \textit{ON} only when the temperature
    exceeds the upper threshold ($\geq 26\,^\circ$C) and turns \textit{OFF} only when
    it drops below the lower threshold ($\leq 24\,^\circ$C). While the temperature is
    between the two thresholds, the alert state remains unchanged. This prevents rapid
    toggling (chattering) when the signal fluctuates near a single threshold value.
  \item \textbf{Debounce filter:} A transition is only accepted after the desired new
    state persists for \texttt{DEBOUNCE\_COUNT} $= 3$ consecutive acquisitions
    (300\,ms). If the condition is not maintained for the full count, the counter resets
    and no state change occurs.
\end{itemize}

Task synchronization is achieved through two FreeRTOS primitives:
\begin{itemize}
  \item A \textbf{binary semaphore} (\texttt{xSemNewData}) given by TaskAcquisition
    after each read and taken by TaskConditioning, ensuring the conditioning logic runs
    exactly once per new sample
  \item A \textbf{mutex} (\texttt{xSensorMutex}) protecting all shared sensor data
    fields against concurrent access by the three tasks
\end{itemize}

\subsection{Hardware and Software Components}

\begin{center}
\begin{tabular}{| p{4cm} | p{10cm} |}
\hline
\textbf{Component} & \textbf{Description \& Role} \\ \hline
\textbf{ESP32 DevKit} & Dual-core microcontroller running FreeRTOS. Hosts all three
  application tasks and provides the OneWire-compatible GPIO and I2C peripherals. \\ \hline
\textbf{DS18B20} & Digital temperature sensor communicating over OneWire protocol on
  GPIO\,4. 9-bit resolution ($\approx 0.5\,^\circ$C), non-blocking conversion
  ($\approx 94$\,ms). \\ \hline
\textbf{Green LED} & Connected to GPIO\,25. Lights up when temperature is within the
  normal range (alert OFF). \\ \hline
\textbf{Red LED} & Connected to GPIO\,26. Lights up when the temperature alert is active
  (temperature confirmed above $26\,^\circ$C). \\ \hline
\textbf{LCD 16$\times$2 I2C} & Connected on the I2C bus (SDA=GPIO\,21, SCL=GPIO\,22),
  address 0x27. Displays real-time temperature and alert state at 500\,ms intervals. \\ \hline
\textbf{$220\,\Omega$ Resistors} & Current-limiting resistors for each LED. \\ \hline
\textbf{Serial-to-USB Interface} & UART at 115200 baud for periodic report output and
  CSV plotter data. \\ \hline
\textbf{FreeRTOS} & Built into the ESP32 Arduino core. Provides preemptive task
  scheduling, binary semaphore, and mutex. \\ \hline
\textbf{PlatformIO} & Build system and IDE integration used for compilation, upload,
  and serial monitoring. \\ \hline
\textbf{Breadboard} & Prototyping base for connecting the DS18B20, LEDs, pull-up resistor,
  and the ESP32 without soldering. \\ \hline
\end{tabular}
\end{center}

\subsection{System Architecture and Solution Justification}

The system is organized into three hardware-abstraction layers following the MCAL /
ECAL / SRV pattern:

\begin{enumerate}
  \item \textbf{MCAL (Microcontroller Abstraction Layer):} Contains only
    \texttt{GpioDriver}, which wraps \texttt{pinMode}, \texttt{digitalWrite},
    \texttt{digitalRead}, and \texttt{analogRead}. All direct register / Arduino
    hardware calls are isolated here.
  \item \textbf{ECAL (Electronic Component Abstraction Layer):} Contains component
    drivers: \texttt{TempSensor} (DS18B20 via OneWire), \texttt{LedDriver},
    \texttt{LcdDisplay} (LiquidCrystal\_I2C), and \texttt{SerialIO}. Each module
    depends only on MCAL or its own library and exposes a clean C API.
  \item \textbf{SRV (Service Layer):} Contains \texttt{SensorData} (shared data store
    with mutex/semaphore), and the three FreeRTOS task implementations:
    \texttt{TaskAcquisition}, \texttt{TaskConditioning}, and \texttt{TaskReporter}.
\end{enumerate}

The data flow proceeds as follows:
\begin{enumerate}
  \item \textbf{TaskAcquisition} reads the DS18B20 sensor every 100\,ms, stores the
    temperature and validity flag in \texttt{SensorData} (under mutex), and gives the
    \texttt{xSemNewData} binary semaphore.
  \item \textbf{TaskConditioning} unblocks immediately on the semaphore, reads the
    latest temperature, applies the \texttt{ProcessChannel} hysteresis + debounce
    function, updates the alert state in \texttt{SensorData}, and drives the green/red
    LEDs accordingly.
  \item \textbf{TaskReporter} runs every 500\,ms via \texttt{vTaskDelayUntil}, reads
    all shared values under mutex, updates the LCD, emits a CSV line for the serial
    plotter, and every 5\,seconds prints a full structured report to STDIO.
\end{enumerate}

\subsection{Relevant Case Study}
A real-world application of hysteresis-based binary signal conditioning is the
\textbf{thermostat control system} found in HVAC equipment. A simple single-threshold
comparator would cause the heating element to switch on and off repeatedly when the
room temperature hovers at exactly the set point, leading to relay wear and unstable
temperature. The hysteresis band (e.g., set point $\pm 1\,^\circ$C) ensures the heater
stays on until the temperature is clearly above the upper threshold and stays off until
clearly below the lower threshold. The debounce timer further prevents brief thermal
noise spikes from triggering spurious transitions -- a technique directly reflected in
the \texttt{DEBOUNCE\_COUNT} parameter of this laboratory work.
